\SourceLeaf{416}%
geraden Linie zugleich angegeben werde, in welchem Sinn sie als beschrieben gedacht wird: ohne eine solche Unterscheidung würde die Richtung einer geraden Linie immer zweien einander entgegengesetzten Radien entsprechen.

\begin{center}2.\end{center}
\noindent\hspace*{1.2em}In dem Hülfskreise nimmt man einen beliebigen Radius als den ersten, oder seinen Endpunkt in der Peripherie als den Anfangspunkt an, und entschliesst sich über den Sinn, in welchem man von diesem aus die Bögen als positiv fortzählen will (ob von der Linken nach der Rechten oder umgekehrt); nach der entgegengesetzten Seite werden die Bögen also als negativ betrachtet. So wird also jede Richtung einer geraden Linie durch Grade etc. oder auch durch eine Zahl ausgedrückt, die dieselben in Theilen des Halbmessers angibt.

Solche Angaben, die um $360^{0}$ oder ein Vielfaches davon verschieden sind, beziehen sich also eigentlich auf gleiche Richtungen, und können, allgemein zu reden, als gleichgültig betrachtet werden; in solchen Fällen jedoch, wo die Entstehungsart einer veränderlichen Neigung mit in Betracht gezogen wird, kann es nöthig sein, dergleichen Angaben sorgfältig zu unterscheiden.

Hat man z. B. sich bestimmt, die Bögen von der Linken zur Rechten zu zählen, und entsprechen zweien geraden Linien $l,l'$ die Richtungen $L,L'$, so ist $L'-L$ die Neigung jener geraden Linien, und man sieht leicht, dass, in so fern $L'-L$ zwischen $-180^{0}$ und $+180^{0}$ fällt, der positive oder negative Werth zugleich anzeigt, dass vom Durchschnittspunkte aus $l'$ rechts oder links von $l$ liegt; allgemein wird diess durch das Zeichen von $\sin(L'-L)$ entschieden.

Es sei $aa'$ ein Stück einer krummen Linie und es entsprechen den Tangenten in $a,a'$ resp. die Richtungen $\alpha,\alpha'$, durch welche Buchstaben nemlich die betreffenden Punkte des Hülfskreises, so wie deren Bogenentfernungen vom Anfangspunkte durch $A,A'$ bezeichnet werden sollen. Die Grösse des Kreisbogens $\alpha\alpha'$ oder $A'-A$ heisst die \textls[100]{Amplitudo} von $aa'$.

Die Vergleichung der Amplitudo des Bogens $aa'$ mit seiner Länge wird uns den Begriff von Krümmung geben. Es sei $l$ indefinite ein Punkt des Bo%
